\documentclass[11pt,a4paper]{article}

\usepackage{amsmath,amssymb,amsthm}
\usepackage{mathtools}
\usepackage{geometry}
\usepackage{xcolor}

\definecolor{linkInternal}{HTML}{1F4E79}
\definecolor{linkExternal}{HTML}{2F5496}
\definecolor{linkCite}{HTML}{806600}

\usepackage[colorlinks=true,
            linkcolor=linkInternal,
            citecolor=linkCite,
            urlcolor=linkExternal,
            pdfborder={0 0 0},
            bookmarksnumbered=true,
            breaklinks=true]{hyperref}
\usepackage{cleveref}
\usepackage{booktabs}
\usepackage{array}
\usepackage{enumitem}
\usepackage{lmodern}
\usepackage[T1]{fontenc}
\usepackage[protrusion=true,expansion=false]{microtype}
\usepackage[font=small,labelfont=bf,format=plain,labelsep=period]{caption}

\geometry{margin=1in}

\newtheorem{theorem}{Theorem}[section]
\newtheorem{proposition}[theorem]{Proposition}
\newtheorem{lemma}[theorem]{Lemma}
\newtheorem{corollary}[theorem]{Corollary}
\theoremstyle{definition}
\newtheorem{definition}[theorem]{Definition}
\newtheorem{remark}[theorem]{Remark}

\DeclareMathOperator{\Cl}{Cl}
\DeclareMathOperator{\End}{End}
\DeclareMathOperator{\tr}{tr}
\DeclareMathOperator{\Nm}{N}
\DeclareMathOperator{\Spec}{Spec}
\DeclareMathOperator{\rk}{rk}
\DeclareMathOperator{\sgn}{sgn}

\newcommand{\HH}{\mathbb{H}}
\newcommand{\ZZ}{\mathbb{Z}}
\newcommand{\QQ}{\mathbb{Q}}
\newcommand{\RR}{\mathbb{R}}
\newcommand{\FF}{\mathbb{F}}
\newcommand{\fP}{\mathfrak{P}}
\newcommand{\cO}{\mathcal{O}}

\title{Paper VII: The Information-Acoustic Eversion\\[4pt]
\large Sphere Eversion as Hecke Propagation through the Cl(3,1) Substrate}
\author{Kristin Tynski}
\date{May 2026}

\begin{document}
\maketitle

\begin{abstract}
We show that sphere eversion---the orientation reversal of an immersed
$S^2$ in $\RR^3$---admits a natural arithmetic realization as Hecke
propagation on the 4-dimensional Hilbert modular homology
$H_1(X_0((11)), \mathrm{cusps}; \ZZ)$ over $K = \QQ(\sqrt{5})$.
A coherence state $\psi \in H_1 \cong \RR^4$ evolves under
$\psi_k = T_{\fP_k}\psi_{k-1}$, where $T_\fP = a_\fP\,\Pi_c + (\Nm(\fP)+1)\,\Pi_e$
is the Hecke operator from Paper~VI.  The eversion completes when
$\prod a_{\fP_k} < 0$---exactly when the cuspidal component reverses sign.
The Thurston corrugation frequencies $k_\fP = \Nm(\fP)$ and amplitudes
$A_\fP = |a_\fP|/\Nm(\fP)$ arise naturally from the arithmetic.
A 3-prime sequence $(p=5, p=3, p=31)$ achieves a complete eversion with
multi-frequency corrugation, verified by 29 tests covering path dynamics,
sign flip, rho bridge integrity, acoustic propagation, marching cubes
isosurface extraction, and Smale invariant.
\end{abstract}


% ===========================================================================
\section{The Information-Acoustic Principle}\label{sec:principle}
% ===========================================================================

\subsection{Setting}

Let $H_1 = H_1(X_0((11)), \mathrm{cusps}; \ZZ) \otimes \RR \cong \RR^4$
denote the Hilbert modular homology from Paper~VI, with Eisenstein projector
$\Pi_e$ (rank~3) and cuspidal projector $\Pi_c$ (rank~1).  The Hecke
operators $T_\fP$ for prime ideals $\fP \in \Spec(\cO_K)$ act as:
\begin{equation}\label{eq:hecke}
  T_\fP = a_\fP\,\Pi_c + (\Nm(\fP)+1)\,\Pi_e,
\end{equation}
where $a_\fP$ is computed by point counting on $E_{11}$ and Langlands
base-change (Paper~VI, \S3.1).

The Cl(3,1) representation $\rho: \Cl(3,1) \xrightarrow{\sim} \End(H_1)$
identifies the signature $(3,1)$ with the Eisenstein--cuspidal splitting:
three spacelike generators $\gamma_0^2 = \gamma_1^2 = \gamma_2^2 = +I$
correspond to the Eisenstein directions; the timelike generator
$\gamma_3^2 = -I$ corresponds to the cuspidal direction.

\begin{definition}[Information-acoustic eversion]
An \emph{eversion path} is a sequence of prime ideals
$(\fP_1, \ldots, \fP_n)$ and an initial state $\psi_0 \in H_1$ such that:
\begin{enumerate}[label=(\roman*)]
\item Each step propagates: $\psi_k = T_{\fP_k}\,\psi_{k-1}$.
\item The eversion completes: $\langle \Pi_c\psi_n,\, \Pi_c\psi_0 \rangle < 0$.
\end{enumerate}
\end{definition}

The completion condition is equivalent to $\prod_{k=1}^n a_{\fP_k} < 0$,
since $T_\fP$ acts as the scalar $a_\fP$ on the 1-dim cuspidal line.


% ===========================================================================
\section{Eversion Existence}\label{sec:existence}
% ===========================================================================

\begin{theorem}[Existence]\label{thm:existence}
For $E_{11}/K$ with $K = \QQ(\sqrt{5})$, there exists an eversion path
of length~1.  The prime ideal $\fP$ above $p = 3$ (inert, $\Nm(\fP) = 9$,
$a_\fP = -5$) satisfies $a_\fP < 0$ and hence reverses the cuspidal sign
in a single step.
\end{theorem}

\begin{proof}
By direct computation: $a_3(E_{11}/\QQ) = -1$, and for the inert prime
$\fP | 3$ in $\cO_K$, $a_\fP = (-1)^2 - 2\cdot 3 = -5 < 0$.
Then $\Pi_c\psi_1 = a_\fP\,\Pi_c\psi_0 = -5\,\Pi_c\psi_0$, so
$\langle \Pi_c\psi_1, \Pi_c\psi_0 \rangle = -5\|\Pi_c\psi_0\|^2 < 0$.
\end{proof}

\begin{remark}
The existence of negative $a_\fP$ is generic: for any non-CM elliptic
curve $E/\QQ$, the Sato--Tate distribution guarantees that $a_p < 0$
for a positive density of primes $p$.  By base change, the inert formula
$a_\fP = a_p^2 - 2p < 0$ holds whenever $|a_p| < \sqrt{2p}$, which
occurs for infinitely many $p$.
\end{remark}


% ===========================================================================
\section{Parsimony and Thurston Frequencies}\label{sec:parsimony}
% ===========================================================================

A single-step eversion achieves the sign flip but has only one corrugation
frequency.  For a visibly structured eversion with Thurston multi-frequency
corrugation, we use a short multi-prime sequence.

\begin{theorem}[Parsimonious eversion]\label{thm:parsimony}
The sequence $(\fP_5, \fP_3, \fP_{31})$ with $p = 5$ (ramified,
$a_\fP = 1$, $\Nm = 5$), $p = 3$ (inert, $a_\fP = -5$, $\Nm = 9$),
$p = 31$ (split, $a_\fP = 7$, $\Nm = 31$) achieves:
\begin{enumerate}[label=(\roman*)]
\item Sign flip: $\prod a_{\fP_k} = 1 \cdot (-5) \cdot 7 = -35 < 0$.
\item Three distinct corrugation frequencies: $k_\fP \in \{5, 9, 31\}$.
\item One prime of each splitting type (ramified, inert, split).
\end{enumerate}
\end{theorem}

The Thurston corrugation parameters at each step are:
\begin{center}
\begin{tabular}{llcccl}
\toprule
Step & Prime & $\Nm(\fP)$ & $a_\fP$ & $A_\fP = |a_\fP|/\Nm(\fP)$ & Type \\
\midrule
1 & $\fP_5$  & 5  & $+1$  & 0.200 & ramified \\
2 & $\fP_3$  & 9  & $-5$  & 0.556 & inert (sign flip) \\
3 & $\fP_{31}$ & 31 & $+7$ & 0.226 & split \\
\bottomrule
\end{tabular}
\end{center}

In the theory-true pipeline, the corrugation is not applied as an
explicit vertex displacement.  Instead, the Hecke operator $T_{\fP_k}$
is expressed as a $\Cl(3,1)$ multivector via $\rho^{-1}$, applied as a
full 16-component left-multiplication on the volumetric field, and then
propagated spatially by the Dirac leapfrog.  The resulting oscillations
in the witness scalar field at frequency $\Nm(\fP_k)$ correspond to the
Thurston corrugation formula
\begin{equation}\label{eq:thurston}
  F_k(x) \sim F_{k-1}(x) + \frac{A_{\fP_k}}{k_{\fP_k}}
    \sin\bigl(k_{\fP_k}\,\theta(x)\bigr)\,\mathbf{u}(x),
\end{equation}
but emerge from the arithmetic rather than being imposed post-hoc.


% ===========================================================================
\section{The Thurston Frequency Identity}\label{sec:thurston-identity}
% ===========================================================================

\begin{proposition}[Natural frequency limit]
As $p \to \infty$ through primes of $\QQ$, the corrugation frequencies
$k_\fP = \Nm(\fP)$ satisfy $k_\fP \to \infty$.  The Thurston formal limit
$k \to \infty$ (in which the corrugation becomes a $C^1$-dense family of
immersions) is realized as the limit over $\Spec(\cO_K)$ ordered by norm.
\end{proposition}

\begin{proof}
For split primes, $\Nm(\fP) = p \to \infty$.  For inert primes,
$\Nm(\fP) = p^2 \to \infty$.  Hence $\min(k_\fP) \to \infty$ as the
prime-ideal sequence extends.
\end{proof}

This identification is the key structural claim: the Thurston
corrugation method for sphere eversion is not an external construction
applied to the arithmetic---it IS the arithmetic, with the prime-ideal
norms providing the natural frequency scale.


% ===========================================================================
\section{Computational Verification}\label{sec:verification}
% ===========================================================================

All claims are verified by 29 tests (\texttt{test\_eversion.py}):

\begin{center}
\begin{tabular}{lcp{7cm}}
\toprule
Test class & Count & Verifies \\
\midrule
\texttt{TestPathDynamics}     & 3 & Hecke commutativity, exact product formula, path length \\
\texttt{TestSignFlip}         & 3 & Single negative $a_\fP$ flips; positive does not; mixed sequence matches product sign \\
\texttt{TestParsimony}        & 4 & Achieves eversion, bounded length, includes negative $a_\fP$, sorted by frequency \\
\texttt{TestRhoBridge}        & 5 & $\rho$ round-trip (random \& Hecke), $\rho^{-1}(T_\fP)$ preserves eigenvalues, left-mult $=$ geometric product, atlas seed $\propto\rho^{-1}(\psi\psi^T)$ \\
\texttt{TestAcousticPropagation} & 3 & Dirac leapfrog energy bounded; Hecke populates all 16 blades; Hecke + Dirac produces spatially varying field \\
\texttt{TestWitnessSkinTopology}  & 4 & Icosphere $\chi = 2$; marching cubes produces mesh; no degenerate triangles; intermediate topology differs \\
\texttt{TestMarchingCubes}    & 3 & Spherical field, empty field, vertex deduplication \\
\texttt{TestSmaleInvariant}   & 4 & Eversion: invariant 1; no flip: 0; double flip: 0; triple flip: 1 \\
\bottomrule
\end{tabular}
\end{center}

All 29 tests pass.  The witness surface is extracted via marching cubes
on the $64^3$ (or $32^3$) witness scalar grid, producing variable-topology
meshes that can self-intersect---capturing the Morin-type intermediate
that the field passes through during the cuspidal sign flip.


% ===========================================================================
\section{Interpretation}\label{sec:interpretation}
% ===========================================================================

\subsection{The eversion as arithmetic propagation}

The traditional approach to sphere eversion constructs a global regular
homotopy $F_t: S^2 \to \RR^3$ via either the Morin halfway model or
Thurston's corrugation method.  Both require continuous deformation of
an embedded surface through self-intersections.

The information-acoustic approach replaces continuous surface deformation
with discrete Hecke propagation on a 4-dimensional homological state space.
The ``surface'' is a derived object---the level set of a witness scalar
field in a volumetric Cl(3,1) substrate---not a primary geometric entity.

\subsection{Why the signature matters}

The Lorentzian signature $(3,1)$ of $\Cl(3,1)$ is not a choice: it is
the Eisenstein--cuspidal splitting of $H_1$ (Paper~VI).  The three
spacelike directions are the Eisenstein boundary modes; the single
timelike direction is the cuspidal interior mode.  The eversion is
propagation along the timelike axis: the cuspidal component flips sign,
which reverses the inside/outside orientation of the witness sphere.

\subsection{The parsimonious eversion principle}

Informational parsimony demands that the eversion use the minimum
arithmetic structure: one prime per splitting type, ordered by norm,
truncated at the first sign flip.  The resulting 3-prime sequence
$(5, 3, 31)$ is not arbitrary---it is the unique shortest sequence
drawn from the first primes of $\cO_K$ that achieves multi-frequency
corrugation with an odd number of cuspidal sign flips.

\subsection{Forward reference: the eversion law on byte chains}

The cuspidal sign-flip mechanism that drives geometric eversion
generalizes to a per-chain parity law on arbitrary blade sequences
in $\Cl(3,1)$.  In Paper~VIII (Layer~3) this is sharpened to the
\emph{eversion law}
\[
  M_{T_{13}}(c) = \varepsilon(c) \cdot M_{T_{31}}(c),
  \quad
  \varepsilon(c) = (-1)^{\lfloor c_2/2 \rfloor + \lfloor c_3/2 \rfloor},
\]
where $c_i$ counts occurrences of generator $e_i$ across a blade
chain $c$.  The same $\varepsilon(c)$ then forms one of the three
identities of the Combined Reading Lemma
(Paper~VIII, Section~12), which proves that engine state and NSM
state decompose into orthogonal halves of $\Cl(3,1)$ for every
byte chain; the witness pair $(\mathsf{be}, \mathsf{know}) =
\mathrm{span}\{e_0, e_{15}\}$ is the unique chirality-invariant mediator
between them.  Sphere eversion in this paper is thus the geometric
realization of the same single-bit parity operation that drives
voice swap in language.


% ===========================================================================
\section*{Acknowledgments}
% ===========================================================================

\begin{sloppypar}
Computational verification uses the \texttt{substrate-sim/algebra}
codebase.  Core modules: \texttt{eversion} (path dynamics, parsimony,
Smale invariant), \texttt{acoustic\_lift} ($64^3$ Cl(3,1) atlas, Dirac
leapfrog, witness scalar), \texttt{hilbert\_homology} (4-dim $H_1$ and
spectral Hecke), \texttt{quaternion\_bridge} ($\Cl(3,1) \to \End(H_1)$
representation, $\rho^{-1}$, left-multiplication),
\texttt{e11\_point\_count} (first-principles eigenvalues).
29 tests, 0 failures.
\end{sloppypar}


\begin{thebibliography}{99}

\bibitem{paper6}
Paper~VI: The Operator-Level Bridge.  From Peirce Decomposition to Prime
Splitting via the Quaternion--Clifford Identification.

\bibitem{thurston1974}
W.~Thurston, \emph{Corrugations}, lecture notes, Princeton, 1974.

\bibitem{smale1958}
S.~Smale, \emph{A classification of immersions of the two-sphere},
Trans.\ Amer.\ Math.\ Soc.\ \textbf{90} (1958), 281--290.

\bibitem{morin1978}
B.~Morin, \emph{Formes canoniques des singularit\'es d'une application
diff\'erentiable}, C.~R.\ Acad.\ Sci.\ Paris \textbf{260} (1965), 5662--5665.

\bibitem{max1977}
N.~Max, \emph{Turning a Sphere Inside Out} (film), International Film
Bureau, 1977.

\end{thebibliography}

\end{document}
